\documentclass[11pt]{article}

\usepackage[margin=1in]{geometry}
\usepackage{amsmath,amssymb,amsthm,mathtools}
\usepackage[T1]{fontenc}
\usepackage{lmodern}
\usepackage{microtype}
\usepackage[colorlinks=true,linkcolor=blue,urlcolor=blue]{hyperref}

\newcommand{\Z}{\mathbf Z}
\newcommand{\F}{\mathbf F}
\newcommand{\ot}{\otimes}
\newcommand{\eps}{\varepsilon}
\newcommand{\Spec}{\operatorname{Spec}}
\newcommand{\id}{\mathrm{id}}
\newcommand{\add}{\mathrm{add}}

\newtheorem{theorem}{Theorem}
\newtheorem{lemma}[theorem]{Lemma}
\theoremstyle{remark}
\newtheorem{remark}[theorem]{Remark}

\title{A Three-Relation Rank-Four Group Scheme\\Not Killed by Four}
\author{}
\date{}

\begin{document}
\maketitle

\begin{abstract}
We give a self-contained rank-four Hopf algebra whose fourth-power word is
not the unit word.  Its base has two generators and only three defining
relations.  These are exactly the two rank-two cancellation relations and
the one mixed relation used by the construction; no Artinian truncation is
added.  Conceptually, the group scheme is a closed subgroup of an elementary
semidirect product \(\mathbb G_a^2\rtimes\mathbb G_m\).  We also prove a
precise, limited minimality statement for this two-parameter recipe.
\end{abstract}

\begin{center}
\fbox{\begin{minipage}{0.92\textwidth}
\small
\textbf{Disclosure: AI-generated work.}\ This document in its entirety---the
construction, the proofs, the computations, and the exposition---was generated
by artificial intelligence. The work was carried out by the AI models Codex
(OpenAI) and Claude (Anthropic).
\end{minipage}}
\end{center}
\medskip

\section{The construction}

Let
\begin{equation}\label{eq:base}
 C=\Z[r,s]/(r^3+2,\ s^3+2,\ rs^2).
\end{equation}
Over \(C\), define
\begin{equation}\label{eq:A}
 A=C[U,V]/(U^2-r^2U+rsV,\ V^2-s^2V),
 \qquad W=UV,
\end{equation}
and put
\begin{equation}\label{eq:lambda}
 \lambda=(1+rU)(1+sV).
\end{equation}

\begin{theorem}\label{thm:main}
The algebra \(A\) is free of rank four over \(C\), on
\[
 1,\qquad U,\qquad V,\qquad W=UV.
\]
It is a Hopf algebra with counit \(\eps(U)=\eps(V)=0\) and coproduct
\begin{equation}\label{eq:Delta}
 \Delta(U)=U\ot1+\lambda\ot U,
 \qquad
 \Delta(V)=V\ot\lambda+1\ot V.
\end{equation}
If \(G=\Spec A\), then
\begin{equation}\label{eq:powers}
 [4]^\#(U)=2sW\ne0,
 \qquad
 [4]^\#(V)=[4]^\#(W)=0,
 \qquad
 [8]=e.
\end{equation}
In particular, \(G\) is a finite locally free group scheme of rank four
which is not killed by four.
\end{theorem}

Here \([n]\) is the power word \(g\mapsto g^n\).  The notation
\([n]^\#\) denotes its comorphism on \(A\).  We do not need \([n]\) to be a
group homomorphism.

\section{Where the three relations come from}

There is one elementary rank-two calculation behind the example.  Suppose
\(T^2=aT\), put \(g=1+bT\), and try
\[
 \Delta(T)=T\ot1+g\ot T.
\]
If \(T_*=T_1+g_1T_2\), then
\begin{equation}\label{eq:rank2}
 T_*^2-aT_*
 =g_1\bigl(2T_1+a(g_1-1)\bigr)T_2
 =g_1(2+ab)T_1T_2.
\end{equation}
Thus the single equation \(ab=-2\) cancels the cross-term and produces the
usual rank-two Hopf pattern.  Indeed,
\(\Delta(g)=g\ot g\), \(g^2=1\), and one may take \(S(T)=T\).

For the rank-four example, use the two pairs
\[
 (a,b)=(r^2,r),
 \qquad
 (a,b)=(s^2,s).
\]
They require
\[
 r^3=-2,\qquad s^3=-2,
\]
which are the first two relations in \eqref{eq:base}.  The new feature is
the correction term \(-rsV\) in
\[
 U^2=r^2U-rsV.
\]
It joins the two rank-two patterns.

One can see the entire bookkeeping in the following five-line table.  Put
\[
 \alpha=r^2,\qquad \beta=-rs,\qquad \gamma=s^2.
\]
The bridge construction asks for
\begin{equation}\label{eq:five-checks}
\begin{array}{c|c}
\text{compatibility identity}&\text{reason in \(C\)}\\ \hline
\alpha r=-2&r^3=-2\\
\gamma s=-2&s^3=-2\\
\gamma r=0&rs^2=0\\
\beta s=0&-rs^2=0\\
\beta r=-\alpha s&-r^2s=-r^2s.
\end{array}
\end{equation}
So the two cubic relations supply the two diagonal cancellations, while
the single relation \(rs^2=0\) kills both unwanted mixed terms.  The bridge
balance is automatic because of the choice \(\beta=-rs\).  This is the
conceptual reason that only three base relations appear.

\section{The one nonzero base element we need}

The relations in \(C\) immediately imply
\begin{equation}\label{eq:base-consequences}
 2r=0,\qquad 2s^2=0,\qquad 4=0.
\end{equation}
Indeed,
\[
 2r=r(s^3+2)-s(rs^2)=0,
\]
and similarly
\[
 2s^2=s^2(r^3+2)-r^2(rs^2)=0.
\]
Multiplying \(2r=0\) by \(r^2\) and using \(r^3=-2\) gives \(4=0\).

What matters is that \(2s\) does \emph{not} vanish.  Here is a short proof
which computes only what we need.  Set
\[
 C_1=\Z[r,s]/(r^3+2,\ s^3+2).
\]
The two cubics are monic, so \(C_1\) is free over \(\Z\) on
\[
 r^is^j,\qquad 0\le i,j\le2.
\]
Define an additive map
\[
 \ell:C_1\longrightarrow\Z/4
\]
by taking the coefficient of the basis element \(s\), modulo \(4\).
We claim that \(\ell\) kills the ideal \((rs^2)\).  Indeed,
\[
 (r^is^j)rs^2=r^{i+1}s^{j+2}.
\]
After using \(r^3=s^3=-2\), this can have an \(s\)-term only when
\(i=j=2\); in that case
\[
 r^3s^4=(-2)(-2s)=4s.
\]
Thus the coefficient of \(s\) is always divisible by \(4\).  The map
\(\ell\) therefore descends to an additive map
\[
 \bar\ell:C\longrightarrow\Z/4
\]
and \(\bar\ell(2s)=2\).  In particular,
\begin{equation}\label{eq:2s-nonzero}
 2s\ne0\quad\text{in }C.
\end{equation}
This is the only structural fact about the base needed below.

\section{Why this is a Hopf algebra}

\subsection*{Freeness and the basic identities}

The relation \(V^2=s^2V\) is monic in \(V\), and after adjoining \(V\), the
relation \(U^2=r^2U-rsV\) is monic in \(U\).  Hence
\[
 A=C\{1,U,V,W\}
\]
is free of rank four.

The multiplication rules needed below are especially sparse:
\begin{equation}\label{eq:W-rules}
 UW=r^2W,\qquad VW=s^2W,\qquad W^2=0.
\end{equation}
Put
\[
 L=1+rU,\qquad M=1+sV,\qquad
 \theta=\lambda-1=rU+sV+rsW.
\]
Using \eqref{eq:base-consequences} and the two quadrics gives
\begin{equation}\label{eq:lambda-identities}
 M^2=1,\qquad
 L^2=1+2sV,\qquad
 \lambda^2=1+2sV,
\end{equation}
and
\begin{equation}\label{eq:theta-identities}
 2\theta=2sV,\qquad
 \theta^2=0,\qquad
 \lambda^{-1}=1-\theta.
\end{equation}
Thus \(\lambda\) is a unit.

\subsection*{The ambient group}

The coproduct formulas in \eqref{eq:Delta} are not isolated guesses.  Let
\[
 H=\Spec B,\qquad B=C[u,v,z,z^{-1}],
\]
where \(H\) has multiplication
\begin{equation}\label{eq:ambient-law}
 (u,v,z)(u',v',z')
 =\bigl(u+zu',\ vz'+v',\ zz'\bigr).
\end{equation}
This law is simply block-diagonal matrix multiplication under
\[
 (u,v,z)\longmapsto
 \begin{pmatrix}z&u\\0&1\end{pmatrix}
 \oplus
 \begin{pmatrix}1&v\\0&z\end{pmatrix}.
\]
In particular, associativity and inverses are automatic.
Equivalently, the coordinate Hopf algebra of \(H\) has
\begin{align*}
 \Delta_H(u)&=u\ot1+z\ot u,\\
 \Delta_H(v)&=v\ot z+1\ot v,\\
 \Delta_H(z)&=z\ot z.
\end{align*}
This is a familiar semidirect product.  Indeed, after replacing \(v\) by
\(w=v/z\), the law becomes
\[
 (u,w,z)(u',w',z')
 =\bigl(u+zu',\ w+z^{-1}w',\ zz'\bigr);
\]
thus \(\mathbb G_m\) acts on the two additive coordinates with weights
\(1\) and \(-1\).

Because \(\lambda\) is a unit, the homomorphism
\[
 B\longrightarrow A,\qquad
 u\longmapsto U,\quad v\longmapsto V,\quad z\longmapsto\lambda
\]
realizes \(G=\Spec A\) as the closed subscheme of \(H\) defined by
\begin{equation}\label{eq:ambient-ideal}
 u^2-r^2u+rsv=0,\qquad
 v^2-s^2v=0,\qquad
 z=(1+ru)(1+sv).
\end{equation}
Therefore \(A\) will be a Hopf algebra once we show that this closed
subscheme contains the identity and is closed under the ambient
multiplication and inverse.  Associativity will then be inherited from
\(H\).  This is the conceptual point of the construction.

\subsection*{Closure under multiplication}

Write \(U_i,V_i,\lambda_i\) for the elements in the two tensor factors, and
let
\[
 U_*=U_1+\lambda_1U_2,
 \qquad
 V_*=V_1\lambda_2+V_2.
\]
By \eqref{eq:ambient-law}, these are exactly the \(u\)- and \(v\)-coordinates
of the product of two points of \(G\).  Four consequences of the three base
relations are
\begin{align}
 s^2(\lambda-1)&=-2V,\label{eq:aux1}\\
 rs(\lambda-1)&=r^2sU,\label{eq:aux2}\\
 rs(\lambda^2-1)&=0,\label{eq:aux3}\\
 2\lambda U+r^2(\lambda^2-\lambda)+r^2sV&=0.\label{eq:aux4}
\end{align}
They make the two quadratic checks short:
\begin{align*}
 V_*^2-s^2V_*
 &=V_1\lambda_2\bigl(s^2(\lambda_2-1)+2V_2\bigr)=0,
\\
 U_*^2-r^2U_*+rsV_*
 &=-rsV_1(1-\lambda_2)
   +\bigl(2\lambda_1U_1+r^2(\lambda_1^2-\lambda_1)\bigr)U_2\\
 &\hspace{2em}-rs(\lambda_1^2-1)V_2\\
 &=\bigl(r^2sV_1+2\lambda_1U_1
          +r^2(\lambda_1^2-\lambda_1)\bigr)U_2=0.
\end{align*}

It remains to check the graph equation in \eqref{eq:ambient-ideal}.  Set
\[
 E=rsV_1U_2.
\]
The factorization has a simple source.  If
\[
 L_*=1+rU_*,\qquad M_*=1+sV_*,
\]
then the ambient product formulas give
\[
 L_*=L_1(1+rM_1U_2),\qquad
 M_*=(1+sV_1L_2)M_2.
\]
Expanding only the two middle factors now gives
\begin{equation}\label{eq:group-like-factor}
 (1+rU_*)(1+sV_*)-\lambda_1\lambda_2
 =L_1M_2E\bigl(2+rU_2+sV_1+E\bigr).
\end{equation}
Every term on the right vanishes.  Indeed,
\[
 2E=E(sV_1)=E^2=0,
\]
and
\[
 E(rU_2)
 =r^2sV_1(r^2U_2-rsV_2)=0
\]
by \(2r=2s^2=0\).  Hence
\[
 \Delta(\lambda)=\lambda\ot\lambda.
\]
Thus the product again satisfies both quadrics and the graph equation.
The identity \((0,0,1)\) plainly lies in \(G\), so ambient multiplication
and identity restrict to \(G\).

\subsection*{Closure under inverses}

The inverse in the ambient group is
\[
 (u,v,z)^{-1}=(-z^{-1}u,\,-vz^{-1},\,z^{-1}).
\]
Restricting this formula to \(G\) suggests
\begin{equation}\label{eq:antipode}
 S(U)=U-sW+r^2sV,\qquad
 S(V)=V+rW.
\end{equation}
These formulas are simply
\[
 S(U)=-\lambda^{-1}U,\qquad
 S(V)=-V\lambda^{-1}.
\]
Using \eqref{eq:W-rules}, direct substitution gives
\[
 S(U)^2-r^2S(U)+rsS(V)=0,\qquad
 S(V)^2-s^2S(V)=0,\qquad
 S(\lambda)=\lambda^{-1}.
\]
Thus \(S\) is a well-defined algebra map.  The convolution identities
follow at once from \eqref{eq:Delta}; for example,
\[
 U+\lambda S(U)=0,\qquad
 S(U)+S(\lambda)U=0,
\]
and similarly for \(V\).  Hence the ambient inverse also restricts to
\(G\).

We have proved that \(G\) is a closed subgroup scheme of \(H\).
Equivalently, the ideal in \eqref{eq:ambient-ideal} is a Hopf ideal in
\(B\), and its quotient is \(A\).  All Hopf axioms for \(A\), including
coassociativity, are therefore inherited from the elementary ambient group.

\section{The fourth-power carry}

We use one elementary norm observation.  If
\[
 \Delta(z)=z\ot1+l\ot z
\]
with \(l\) group-like, then repeated multiplication gives
\[
 [n]^\#(z)=N_n(l)z,
 \qquad
 N_n(l)=1+l+\cdots+l^{n-1}.
\]
For a right skew-primitive element, the formula is
\([n]^\#(z)=zN_n(l)\).

Since \(\lambda=1+\theta\) and \(\theta^2=0\),
\[
 N_4(\lambda)
 =1+\lambda+\lambda^2+\lambda^3
 =4+6\theta
 =2\theta
 =2sV.
\]
Therefore
\[
 [4]^\#(U)=N_4(\lambda)U=2sUV=2sW,
\]
which is nonzero by \eqref{eq:2s-nonzero} and freeness of \(A\).  On the
other hand,
\[
 [4]^\#(V)=VN_4(\lambda)
 =2sV^2=2s^3V=-4V=0.
\]
Since the comorphism of the power word is an algebra map,
\([4]^\#(W)=0\) as well.

Finally, \(\lambda^4=1\), so
\[
 N_8(\lambda)
 =N_4(\lambda)(1+\lambda^4)
 =2N_4(\lambda)=0.
\]
Both generators \(U,V\) go to zero under \([8]^\#\), and hence
\([8]=e\).  This proves Theorem~\ref{thm:main}.

\section{In what sense three relations are minimal}

The number of equations in an arbitrary ring presentation is not an
intrinsic invariant, so the appropriate claim is modest: the particular
three-relation compatibility locus used above cannot itself be defined
inside \(\Spec\Z[r,s]\) by fewer equations.

Let
\[
 I=(r^3+2,\ s^3+2,\ rs^2)\subset\Z[r,s],
 \qquad
 \mathfrak m=(2,r,s).
\]
We have already seen that \(2r\in I\), and it follows that
\[
 4\in I,\qquad r^4\in I,\qquad s^6\in I.
\]
Thus \(\sqrt I=\mathfrak m\).  Since \(\mathfrak m\) has height three,
Krull's height theorem says that \(I\) cannot be generated by fewer than
three elements.  More generally, if \(J\supset I\) is any proper ideal,
then \(\sqrt J=\mathfrak m\), so \(J\) also needs at least three generators.
Thus this particular compatibility locus cannot be cut out inside
\(\Spec\Z[r,s]\) by fewer than three equations.

This is not a claim about every conceivable construction of a rank-four
Hopf algebra.  It says exactly what is useful here: after fixing the two
parameters \(r,s\) and the coefficients
\[
 \alpha=r^2,\qquad \beta=-rs,\qquad \gamma=s^2,
\]
no fourth relation, nilpotence cutoff, or Artinian truncation is being
smuggled into the proof.  The three base relations themselves
happen to force the base to be Artinian, but that is a consequence rather
than an input.

\end{document}
